%!TEX root = TQFTnotes.tex
%%%%%%%%%%%%%%%%%%%%%%%%%%%%%%%%%%%%%%%%%%%%%%%%%%%%%%%%%%
\chapter{Physical motivation}\label{c:physical-motivation}
Overview of classical and quantum field theory. $N$ is the configuration space
of the theory, $\dim N = n-1$; $M= N\times[t_0, t_1]$ is the spacetime.

%%%%%%%%%%%%%%
\section{Classical/quantum mechanics}\label{s:mechanics}

The table below summarizes the main structures in classical
and quantum mechanics:

\begin{table}[h]
  \centering
  \begin{tabular}{r|l|l}
    &Classical mechanics & Quantum mechanics  \\
    \hline
    Position
    &$x\in N$ (configuration space)
    & $\psi \in H_N=L^2(N)$ \\
    Trajectory
    &$\gamma\colon [t_0, t_1]\to N$
    &$U_{t_0,t_1}\colon H_N\to H_N$ \\
    Action
    & $S[\gamma]=\int L(\gamma,\dot\gamma) dt$
    & \\
    Laws of motion
    &\parbox[t]{5cm}{Least action principle\par
      \quad $\delta S[\gamma]=0$ (fixed endpoints)
    }
    & \parbox[t]{7cm}{
      $(U\psi)(x_1) = \int_N (K(x_0, x_1) \psi(x_0) dx_0$\par
        $K(x_0,x_1) = \int_P e^{\frac{i}{\hbar} S[\gamma]} D\gamma$\par
        $P=\{\gamma: \gamma(t_0)=x_0, \gamma(t_1)=x_1$\}

      }
      \\
  \end{tabular}
  \caption{FIXME}
  \label{tab:mechanics}
\end{table}
Note that in the limit $\hbar \to 0$, the integral in the definition of
$K(x_0, x_1)$ is dominated by the contribution from the neighborhood of
critical points of the action. For finite-dimensional integrals of this
form, method of stationary phase allows one to write asymptotic expansion
in powers of $\hbar$.

%%%%%%%%%%%%%%
\section{Classical/quantum field theory}\label{s:field-theory}
In field theory, a point $x\in N$ is replaced by the field $\ph$, which is
typically either a function on $N$ (with values in some target space $V$)
or a section of a vector bundle or a connection. We will denote by
$\mathcal{F}_N$ the space of all fields on $N$.

\begin{table}[h]
  \centering
  \begin{tabular}{r|l|l}
    &Classical field theory  & QFT  \\
    \hline
    &$\ph \in \calF_N$
    & $\psi \in \calH_N=L^2(\calF_N)$ \\
    Trajectory
    &$\ph\in \calF_{M}$, $M=N\times[t_0, t_1]$
    &$U_{t_0,t_1}\colon \calH_N\to \calH_N$ \\
    Action
    & $S[\ph]=\int_M L(\ph, \del_\mu\ph) d^{n}x$
    & \\
    Laws of motion
    &\parbox[t]{5cm}{Least action principle\par
      \quad $\delta S[\ph]=0$ (fixed restriction to $t=t_0$,$t=t_1$)
    }
    & \parbox[t]{7cm}{
      $(U\psi)(\ph_1) = \int_N (K(\ph_0, \ph_1) \psi(\ph_0) D\ph_0$\par
        $K(\ph_0,\ph_1) = \int_P e^{\frac{i}{\hbar} S[\ph]} D\ph$\par
        $P=\{\ph\in \calF_M: \ph(t_0)=\ph_0, \ph(t_1)=\ph_1$\}
      }\\
  \end{tabular}
  \caption{FIXME}
  \label{tab:qft}
\end{table}
The above formulas freely use integrals over infinite-dimensional
spaces of paths/fields, which are very hard to define rigorously; as of
now, QFT cannot be defined rigorously as a mathematical theory except very
few special cases. Usually the best one can do is to write the
asymptotics formulas, expanding the operator $U$ or the kernel $K$ as a
series in powers of $\hbar$, and even that requires a lot of work.

However, if we ignore that and assume that all the above makes sense, we get the
following structure (again, at the moment we choose to ignore details)
\begin{enumerate}
  \item For every $n-1$-dimensional manifold $N$,
    we have a Hilbert space $\calH_N$. Moreover, since
    $\calF_{N_1\sqcup N_2} = \calF_{N_1}\times \calF_{N_2}$, we get
    \begin{equation}\label{e:qft-axiom1}
      \calH_{N_1\sqcup N_2} = \calH_{N_1}\otimes \calH_{N_2}
    \end{equation}
  \item The definition of operator $U_{t_0, t_1}$ immediately generalizes,
    using the same formulas, to arbitrary $n$-dimensional manifold $M$ and
    not only the cylinder $N\times [t_0, t_1]$: if $\del M= N_0\sqcup N_1$
    (with appropriate orientations), then $U_M$ is an operator
    \begin{equation}
      U_M\colon \calH_{N_0}\to \calH_{N_1}
    \end{equation}
    In particular, for a closed spacetime (i.e. $\del M=\varnothing$), we get
    $U_M\in \C$ is a number; in physics it is usually called the
    \emph{partition function} and denoted by $Z(M)$.
  \item So defined operators satisfy the gluing axiom:
    (FIXME)
\end{enumerate}

So far we haven't specified what we mean by ``manifold''. Usually in
physics the spacetime is an oriented pseudo-Riemannian manifold, with metric
of signature $(1, n-1)$. However, we will be interested in a very special
case, namely when all the structures above do not use the metric on the
space; thus, we only require $M$ to be an oriented smooth manifold. Such
theories are called \emph{topological}. The most famous example of such a
theory is Chern--Simons theory (see \cite{witten-jones}), in which the
spacetime $M$ must be 3-dimensional, fields are $\SU(n)$ connections on $M$
(which,
after choosing a trivialization of the bundle, can be described by one-forms
with values in the Lie algebra  $\su(n)$) and the action is given by
\begin{equation}
  S[A] = g \int_M \tr(A\wedge dA + \tfrac{2}{3}A\wedge A\wedge A)
\end{equation}

We can take properties above as axioms.
\begin{predefinition}
  An $n$-dimensional topological field theory (TQFT) is the following
  collection of data
  \begin{itemize}
    \item For every  $(n-1)$-dimensional oriented smooth manifold $N$, a
      vector space $Z(N)$
    \item For every $n$-dimensional oriented manifold $M$ such that
      $\del M=N_0\sqcup N_1$, a linear operator $Z_M\colon Z_{N_0}\to Z_{N_1}$,
  \end{itemize}
  This data has to satisfy the properties above (gluing axiom, disjoint
  union).
\end{predefinition}

(Clearly, there is a lot more detail to be filled in, which we will do later).

This definition is due to Atiyah (\cite{atiyah}).

Having this definition, instead of defining a TQFT by a path integral or other
global construction (which can be very hard to do), we can try to construct
TQFTs by defining $Z(M)$ for ``elementary blocks'' and gluing every manifold
from them. For $n=2$, these blocks are disk, cylinder, and pair of pants.

\begin{center}
  \begin{tikzpicture}
    \pic at (0,0) {pants};
    \pic at (2,0) {cylinder};
    \pic at (3, -0.5) {cup};
  \end{tikzpicture}
\end{center}

Of course, we need to have some compatibility conditions to ensure that
different ways of gluing the manifold from pieces give the same result;
essentially, we are trying to define manifolds by ``generators and
relations''. For $n=2$ it is not too hard (we will do so next week); for
larger $n$, it is harder --- e.g. for $n=3$ we can consider so-called
Heegard decomposition (which already requires infinitely many blocks, all
handlebodies of arbitrary genus) or use  surgery along knots
(again, to be discussed later).


However, there is another idea: why stop at manifolds with boundary?
If we allow manifold with corners of any codimension, the set of
``elementary blocks'' is greatly simplified: any manifold can be glued
together from just one block, an $n$-dimensional simplex.
Such theories are called  fully extended (or sometimes local) TQFTs;
see \cite{lurie-cobordism}.

But the price you pay is algebraic complexity: for closed $n$-manifold,
the theory gives you a number;  for a codimension 1 manifold, a vector space.
What is the next level?

One way to answer is that codimension 2 should give you a category, and so on;
thus, doing it properly requires a notion of $n$-categories (or a variation of
that, called $(\infty, n)$ categories). This gets confusing quite fast\dots

\clearpage
